\section{CCM architecture}

\subsection{System overview}

Figure~\ref{fig:ccm-architecture} is the canonical architecture. A task or environment supplies observations and available actions to an orchestrator. Each column has a local model, one reference-frame state, local observations, candidate hypotheses, and a specialization label. Columns exchange structured states rather than private chain-of-thought. An append-only memory records observations with provenance. Consensus either commits a compatible hypothesis, abstains on low confidence, or remains unresolved. If the evidence is insufficient and an action is available, the active policy chooses a discriminative observation and the loop continues.

\begin{figure}[t]
\centering
\begin{tikzpicture}[scale=0.52, transform shape,
  node distance=7mm and 10mm,
  box/.style={draw, rounded corners, align=center, minimum height=9mm, minimum width=25mm, fill=gray!8},
  state/.style={draw, rounded corners, align=center, minimum height=11mm, minimum width=31mm, fill=blue!8},
  decision/.style={draw, diamond, aspect=2, align=center, fill=orange!15},
  link/.style={-{Latex[length=2mm]}, thick}
]
\node[box] (task) {Task / environment};
\node[state, right=of task] (orch) {Orchestrator\\budget and trace};
\node[box, above right=of orch] (c1) {Column 1\\local model + frame};
\node[box, right=of orch] (c2) {Column 2\\local model + frame};
\node[box, below right=of orch] (c3) {Column 3\\local model + frame};
\node[state, right=25mm of c2] (mem) {Structured\\append-only memory};
\node[state, above=of mem] (hyp) {Distributed\\hypotheses and votes};
\node[decision, right=of mem] (vote) {compatible\\quorum?};
\node[box, below right=of vote] (active) {Disambiguating\\action / tool};
\node[box, above right=of vote] (out) {Commit, abstain,\\or unresolved};
\draw[link] (task) -- (orch);
\draw[link] (orch) -- (c1);
\draw[link] (orch) -- (c2);
\draw[link] (orch) -- (c3);
\draw[link] (c1) -- (mem);
\draw[link] (c2) -- (mem);
\draw[link] (c3) -- (mem);
\draw[link] (c1) -- (hyp);
\draw[link] (c2) -- (hyp);
\draw[link] (c3) -- (hyp);
\draw[link] (hyp) -- (vote);
\draw[link] (mem) -- (vote);
\draw[link] (vote) -- node[above, sloped] {yes} (out);
\draw[link] (vote) -- node[below, sloped] {no / insufficient evidence} (active);
\draw[link] (active) -| (task);
\end{tikzpicture}
\caption{CCM architecture. The arrows represent explicit state and observation contracts; the diagram does not imply biological correspondence.}
\label{fig:ccm-architecture}
\end{figure}

\subsection{Reference frames as computation}

A CCM reference frame is not a persona, topic, namespace, domain, or agent perspective. It is a navigable state space
\[
  F=(f, k, o, s, T, O),
\]
where $f$ is a stable identity, $k$ is the frame kind, $o$ is the origin, $s$ is the current state, $T$ is a set of validated transitions, and $O$ maps locations to observation identifiers. A transition $(a,s_i,s_j)\in T$ is valid only when the current state is $s_i$; an observation is attachable only when its frame and location are valid.

The first concrete example is spatial/object-like. A sensor observes a mug feature at object-local location \texttt{handle\_left}; action \texttt{rotate\_90} moves the state to \texttt{handle\_top}; the frame predicts an observation at that location. The second is conceptual. A claim begins at \texttt{claim\_unresolved}; action \texttt{inspect\_source\_a} moves to \texttt{claim\_supported} when the source observation supports the claim; a competing claim remains open if the observation does not discriminate. Both examples contain a state, a transition, a predicted next location, and evidence associated with a location.

\subsection{Columns}

A column is a local model plus a bounded state owner. It is not merely one agent or one prompt. Its minimum state contains local observation identifiers, a reference-frame state, one or more candidate hypotheses, evidence identifiers, confidence, uncertainty, provenance, requested next action, and received peer states. Columns may share a base model or use heterogeneous models; v0.1.0 exposes specialization labels but does not claim a trained specialization.

The standard interface is:

\begin{verbatim}
observe(observation)
update_hypotheses(candidates)
propose_action(available_actions)
emit_state()
receive_peer_state(peer_state)
vote()
commit(status)
\end{verbatim}

The interface separates local update from communication. A peer receives a serializable state, not hidden reasoning. A column can be committed, abstained, or unresolved by the orchestrator.

\subsection{CCM memory model}

CCM memory is neither conversation history, a scratchpad, ordinary RAG, a vector index, nor an assertion that episodic, semantic, and procedural memory have been biologically reconstructed. It is a structured, append-only observation ledger. The v0.1.0 object schema is:

\begin{longtable}{@{}p{0.22\linewidth}p{0.69\linewidth}@{}}
\toprule
Field & Contract \\
\midrule
object\_id & Stable unique identity; duplicate writes are rejected. \\
content & Human-readable observation payload. \\
frame\_id, location & Reference-frame identity and state/location. \\
source, provenance & Origin and trace link; required for every object. \\
sequence, timestamp & Ordered run-local history; timestamps are stored at run level. \\
confidence, validation\_state & Explicit quality and external validation status. \\
relations & Structured relations such as support or contradiction. \\
embedding & Optional future field; not used by the v0.1.0 symbolic adapter. \\
\bottomrule
\end{longtable}

Historical observations cannot be silently mutated. A model-generated hypothesis is not accepted as an external observation. Conflicting observations remain separate records and are visible to consensus and the run trace.

\subsection{Structured voting}

A vote is a tuple containing \texttt{vote\_id}, \texttt{column\_id}, \texttt{hypothesis\_id}, \texttt{hypothesis}, and confidence. A hypothesis contains a claim, frame, location, evidence identifiers, provenance, confidence, uncertainty, and optional requested action. Two hypotheses are compatible only when their normalized claim, frame identity, and location agree. Structured consensus sums confidence within compatible groups, requires a minimum quorum and confidence, and requires a margin over the runner-up. Ties and low-confidence groups are not forced into an answer. This differs from majority voting, self-consistency, debate, confidence-weighted answer aggregation, and judge-model aggregation because the primary unit is a frame-associated hypothesis with explicit evidence.

\subsection{Active evidence}

The active loop maintains competing hypotheses, identifies missing support, scores available actions by expected information gain divided by cost, executes one valid action, appends the resulting observation, and redistributes state. An action is not an answer and cannot create evidence without an environment response. If the action budget is exhausted, the system terminates with the best permitted non-commit outcome. The harness records action selection, returned observation, cost, and whether required evidence was retrieved.

\subsection{Primary algorithm}

Algorithm~\ref{alg:ccm} is implementation-independent pseudocode. It makes termination and disagreement explicit.

\begin{enumerate}
  \item Initialize $N$ columns, their reference frames, an empty append-only memory, and a run manifest.
  \item Observe the task/environment and write only externally returned observations with source and provenance.
  \item Each column updates local hypotheses and emits a structured state.
  \item Exchange states and group compatible hypotheses by frame, location, and claim.
  \item If a group satisfies quorum, confidence, and margin criteria, commit it.
  \item Otherwise, if uncertainty remains and an action budget is available, choose the highest information-gain-per-cost action.
  \item Execute the action, validate its returned observation, append it, and repeat from step 3.
  \item If evidence is insufficient, confidence is below threshold, or the budget is exhausted, abstain or terminate unresolved.
\end{enumerate}

\begin{table}[t]
\centering
\caption{CCM termination outcomes.}
\label{alg:ccm}
\begin{tabular}{ll}
\toprule
Outcome & Condition \\
\midrule
Committed & Compatible structured quorum, confidence, and margin are met. \\
Abstained & No safe commitment; aggregate confidence is below threshold. \\
Unresolved & Competing hypotheses remain too close or lack quorum. \\
\bottomrule
\end{tabular}
\end{table}

\subsection{Invariants}

The implementation and tests enforce the following invariants where feasible: every memory object has provenance; every hypothesis has a reference frame; every vote references an existing hypothesis; historical observations cannot be overwritten; disagreement is retained with its resolution mechanism; model-generated information cannot be represented as externally observed evidence; and each run has a deterministic configuration identity plus model, task, hardware, seed, and git metadata.
